\section{Why Develop Software in Go?}
Go, also known as Golang, was developed by Google and first released publicly in 2009. It was designed to address challenges associated with large-scale software development, particularly long compilation times, complex dependency management, and the difficulty of writing concurrent systems in some established languages. Since its introduction, Go has evolved into a mature general-purpose programming language with improvements to its compiler, tooling, standard library, and language features. A significant development was the introduction of generics in Go 1.18, which provided greater flexibility for reusable code while retaining Go's emphasis on simplicity.\\

Go is particularly well suited for modern systems and backend development because it combines relatively simple syntax with strong static typing, efficient compilation, and built-in support for concurrency through goroutines and channels. Its standard toolchain also provides integrated facilities for formatting, testing, dependency management, and building applications, reducing the need for extensive external tooling. Go produces compiled binaries and has efficient runtime performance, making it suitable for services, cloud infrastructure, networking applications, and distributed systems.\\

Developing in Go can therefore provide a balance between developer productivity and system performance. Its deliberately small language design makes code relatively easy to read and maintain, while its concurrency model and cross-platform compilation support modern distributed applications. Although Go is not the best choice for all types of software, its simplicity, strong tooling, performance, and suitability for scalable systems make it a practical choice for developing reliable and maintainable applications.